\documentclass[12pt,a4paper]{article}
\usepackage[utf8]{inputenc}
\usepackage[T1]{fontenc}
\usepackage[french]{babel}
\usepackage{hyperref}
\usepackage{graphicx}
\usepackage{geometry}
\usepackage{tcolorbox}
\usepackage{enumitem}

\geometry{top=2.5cm, bottom=2.5cm, left=2.5cm, right=2.5cm}

\title{\textbf{Guide Étudiant : Procédure de Mobilité Internationale}\\ \large Manuel Pratique d'Utilisation de la Plateforme}

\begin{document}

\maketitle
\tableofcontents
\newpage

\section{Introduction et Accès}
Bienvenue sur le module de mobilité internationale. Ce guide vous expliquera pas à pas comment consulter les universités partenaires, formuler vos vœux, et répondre à votre proposition d'affectation.

Pour accéder à la plateforme de mobilité :
\begin{tcolorbox}[colback=blue!5!white,colframe=blue!75!black,title=Accès au module]
Connectez-vous à votre espace étudiant, puis cliquez sur l'onglet \textbf{« Carte de Mobilité »} situé dans la barre de navigation principale.
\end{tcolorbox}

\section{Consultation et Filtres du Catalogue}
Une fois sur la page de mobilité, vous ferez face à une carte interactive présentant les universités partenaires. 

\begin{tcolorbox}[colback=blue!5!white,colframe=blue!75!black,title=Affichage Automatique par Filière]
\textbf{Information importante :} Le système trie automatiquement les universités pour vous. La carte \textbf{n'affiche que les universités qui proposent des places pour votre filière}. Vous n'avez donc pas besoin de chercher à filtrer par filière, tout ce qui apparaît à l'écran correspond exactement à ce à quoi vous avez le droit de postuler !
\end{tcolorbox}

\begin{itemize}[leftmargin=*]
    \item \textbf{Exploration :} Vous pouvez cliquer sur les différents marqueurs sur la carte pour découvrir les universités, ou parcourir le tableau situé juste en dessous.
    \item \textbf{Filtres :} Utilisez le menu déroulant pour filtrer les universités uniquement par \textbf{Semestre (S8 ou S9)}.
    \item \textbf{Informations détaillées :} Pour chaque université, vous pouvez voir le nombre de places disponibles, les langues requises, ainsi qu'un lien vers leur site web officiel.
\end{itemize}

\section{Formulation de vos Vœux}
C'est l'étape la plus importante de votre procédure. Vous devez sélectionner les universités où vous aimeriez aller, et les classer par ordre de préférence.

\subsection{Ajout d'un Vœu}
Pour ajouter une université à votre liste de vœux :
\begin{enumerate}
    \item Cherchez l'université de votre choix directement sur la \textbf{carte des vœux}.
    \item Cliquez sur l'université, puis sur le bouton \textbf{« Ajouter aux voeux »} qui s'affiche dans le panneau d'information.
    \item L'université apparaîtra immédiatement dans votre \textbf{Panier de vœux}.
\end{enumerate}

\begin{tcolorbox}[colback=yellow!10!white,colframe=orange!80!black,title=Règles sur la quantité de vœux]
\begin{itemize}
    \item Vous pouvez faire \textbf{jusqu'à 5 vœux maximum}.
    \item Il n'est \textbf{pas obligatoire d'atteindre ce maximum}. Vous pouvez tout à fait soumettre votre dossier avec seulement 1, 2, 3 ou 4 vœux si vous ne souhaitez aller nulle part ailleurs.
\end{itemize}
\end{tcolorbox}

\subsection{Cas Particulier : Le Choix du Stage}
Si vous ne souhaitez pas effectuer de mobilité universitaire classique mais que vous désirez faire un \textbf{stage en entreprise} à l'étranger, vous devez tout de même l'indiquer sur la plateforme.

Le système utilise un "faux choix" d'université pour désigner le stage :
\begin{itemize}
    \item Recherchez l'établissement nommé \textbf{« École Polytech Annecy-Chambéry »} directement sur la carte des vœux.
    \item \textbf{Attention au filtre :} Ce faux choix n'apparaît que si votre filtre de semestre est réglé sur \textbf{« Tout »} ou sur \textbf{« S8 »}. Il n'apparaîtra pas si vous filtrez par « S9 ».
    \item \textbf{Signification :} Même s'il apparaît sous le filtre S8, sélectionner ce choix signifie simplement que vous optez pour la mobilité \textbf{Stage} en général, et non pas que votre stage est strictement limité au S8.
\end{itemize}

\begin{tcolorbox}[colback=red!5!white,colframe=red!75!black,title=⚠️ OBLIGATION DE SOUMISSION]
Même si aucune des universités partenaires ne vous intéresse, \textbf{vous devez impérativement ajouter ce faux choix "École Polytech Annecy-Chambéry"} à votre panier et soumettre votre dossier pour officialiser votre volonté de faire un stage.
\end{tcolorbox}

\subsection{Ordonnancement : L'importance de la Priorité}
L'ordre dans lequel vous classez vos vœux dans votre panier est \textbf{absolument critique}. L'algorithme d'affectation lira toujours votre liste de haut en bas. 

\begin{tcolorbox}[colback=red!5!white,colframe=red!75!black,title=⚠️ STRATÉGIE DE CLASSEMENT]
Votre \textbf{meilleur choix (votre rêve absolu) doit impérativement être en position Numéro 1} tout en haut de la liste.
\newline \newline L'algorithme essaiera toujours de vous donner votre vœu N°1 d'abord. S'il n'y a plus de place, il passera à votre vœu N°2, et ainsi de suite.
\end{tcolorbox}

Pour modifier l'ordre :
\begin{itemize}
    \item Utilisez les flèches \textbf{« Monter »} ($\uparrow$) et \textbf{« Descendre »} ($\downarrow$) à côté de chaque université dans votre panier pour ajuster précisément votre classement.
    \item Vous pouvez également supprimer une université de votre panier en cliquant sur l'icône de la corbeille.
\end{itemize}

\subsection{Soumission Définitive}
Tant que vous n'avez pas soumis votre dossier, il est considéré comme un simple brouillon.
\begin{enumerate}
    \item Une fois que votre liste est parfaite et classée dans le bon ordre, cliquez sur le bouton \textbf{« Soumettre mes vœux »}.
    \item \textbf{Attention (Panier verrouillé) :} Une fois soumis, votre panier est verrouillé. Si vous avez fait une erreur de clic ou un mauvais choix de dernière minute, \textbf{contactez très rapidement le service RI (avant la date limite)}. Ils ont un bouton de réinitialisation sur leur interface qui leur permet d'annuler votre soumission et de vous rendre la main.
\end{enumerate}

\begin{tcolorbox}[colback=red!5!white,colframe=red!75!black,title=⚠️ CONSÉQUENCE DU RETARD DE SOUMISSION]
Si vous oubliez de cliquer sur "Soumettre" et que la date limite officielle fixée par l'administration est franchie, le système verrouillera votre panier dans son état actuel et \textbf{le soumettra automatiquement de force} à votre place. 
\newline Si votre panier était vide à ce moment-là, vous serez exclu de la procédure et l'algorithme ne vous affectera nulle part. Soyez responsables !
\end{tcolorbox}

\section{Résultats et Décision d'Affectation}
Une fois la période de soumission terminée, l'administration lancera l'algorithme d'affectation automatique.

Lorsque l'algorithme aura tourné, vous serez informés directement par le \textbf{Service des Relations Internationales (RI)} qu'il est temps de vous connecter pour consulter vos résultats. Le système lui-même n'envoie pas de notification automatique. Une fois connecté, un nouveau panneau s'affichera sur votre module de mobilité pour vous indiquer l'université qui vous a été attribuée.

\subsection{Cas de Non-Affectation}
Il est possible que l'algorithme n'ait pas pu vous affecter à l'un de vos vœux (par exemple, si vous avez fait peu de choix et que toutes les places étaient déjà prises par des étudiants mieux classés). 
\newline Dans ce cas précis :
\begin{itemize}
    \item \textbf{Vous ne recevrez aucune notification du système.}
    \item Votre interface ne changera pas par rapport à la phase de soumission.
    \item Vous n'aurez aucun bouton pour "Accepter" ou "Refuser" quoi que ce soit.
\end{itemize}
Si vous constatez cela, attendez simplement que le Service RI vous contacte directement pour envisager avec vous les suites à donner à votre dossier.

\subsection{Accepter ou Refuser (Si vous êtes affecté)}
Face à cette proposition, vous avez un délai limité pour donner votre réponse ferme :
\begin{enumerate}
    \item \textbf{Accepter :} Cliquez sur le bouton vert \textbf{« Accepter »}. L'affectation est alors définitivement validée, et vous pourrez entamer vos démarches administratives.
    \item \textbf{Refuser :} Si, pour une raison personnelle, vous décidez de ne plus partir dans cette université, cliquez sur \textbf{« Refuser »}. 
\end{enumerate}

\begin{tcolorbox}[colback=red!5!white,colframe=red!75!black,title=⚠️ ATTENTION AU DÉLAI DE RÉPONSE]
Votre statut sera initialement affiché comme \textit{« En attente »} de votre réponse. Si vous ne cliquez sur aucun bouton avant la date de clôture imposée par l'administration, votre silence sera considéré comme un \textbf{REFUS AUTOMATIQUE} et votre place sera définitivement perdue et réattribuée.
\end{tcolorbox}

\end{document}
